\section{Related Works}
\subsection{Existing Systematic Literature Reviews and Surveys}
The project website identifies five especially relevant prior reviews. The closest methodological competitor is the PRISMA-based SLR on small language models available on the project website. It is the nearest precursor because it combines a transparent review process with a broad overview of SLM research. However, it remains broader than the present study: its main focus is the SLM landscape itself, whereas our review centers on the relationship between SLMs, LLMs, and hybrid architectures.

The remaining relevant works are mostly narrative or semi-systematic surveys. \emph{A Comprehensive Survey of Small Language Models in the Era of LLMs} provides wide coverage of definitions, training methods, applications, collaboration, and trustworthiness, but it does not follow a reproducible selection protocol. \emph{A Survey on Collaborative Mechanisms Between Large and Small Language Models} offers a technically valuable taxonomy of routing, auxiliary, fusion, and distillation strategies, especially for edge deployment, but its scope is mechanism-centered rather than corpus-centered. \emph{A Survey on Collaborating Small and Large Language Models for Performance, Cost-effectiveness, Cloud-edge Privacy, and Trustworthiness} overlaps substantially with the current review's themes, yet its methodology is still narrative rather than PRISMA-guided. Finally, \emph{Collaborative Inference and Learning between Edge SLMs and Cloud LLMs} provides systematic organization of post-2023 edge-cloud work, but it is intentionally narrow in deployment context.

What emerges from these prior reviews is not a lack of attention to SLMs or collaboration, but a difference in analytical emphasis. Previous surveys often answer the question ``What kinds of SLMs or collaboration mechanisms exist?'' The present review instead asks ``What patterns actually recur across recent primary studies, how do they affect performance and efficiency, and which domains benefit most from hybridization?'' This difference is important because the research problem has shifted from model cataloguing to system design.

\subsection{Main Gaps in the Existing Review Literature}
The review website's comparative analysis surfaces four recurring gaps, all of which are visible in the primary studies as well.

First, application domains are underemphasized relative to architecture taxonomies. Prior surveys are strong on categories such as routing, distillation, and cloud-edge cooperation, but weaker on a systematic comparison of concrete domains such as telecom question answering, radiology reporting, optimisation modelling, dialogue state tracking, or industrial video captioning. This matters because whether an SLM is ``good enough'' depends heavily on task structure and domain constraints.

Second, most related works are narrative surveys rather than transparent systematic reviews. Narrative surveys are useful for breadth and conceptual framing, but they make it harder to reproduce study selection, audit exclusions, and compare cross-study evidence using a consistent extraction schema.

Third, prior reviews often stop before the newest wave of 2025 and early-2026 papers that focus on on-device agents, uncertainty-aware inference, prompt optimisation, multimodal collaboration, and energy-aware routing. Because this area evolves rapidly, even a one-year gap meaningfully changes the architectural picture.

Fourth, quantitative synthesis remains thin. Prior surveys usually discuss performance, efficiency, and privacy conceptually, but they rarely align cross-study indicators such as latency in milliseconds, API cost, FLOP reduction, or energy-per-token. This limits their practical usefulness for system builders who need to make deployment decisions rather than only understand the research landscape.

\subsection{Positioning of This Study}
This review is positioned as a methodologically explicit and engineering-oriented synthesis of SLM--LLM systems. It differs from prior work in three ways. First, it treats domain suitability as a first-class analytical axis rather than a side effect of architectural categories. Second, it explicitly maps research questions to extraction fields so that claims about performance, efficiency, and orchestration remain traceable to the underlying evidence. Third, it combines broad corpus-level analysis with representative deep dives into primary papers, which is necessary because the dataset's ``key findings'' fields are valuable for synthesis but too compressed to fully explain why a design works.

The website's own framing is also instructive: it argues that the review should be seen as a PRISMA-guided study of domain-specific applications of SLM--LLM hybrid architectures that emphasizes quantitative performance metrics. Our reanalysis supports this framing. In the current corpus, the most consequential papers are not simply ``about small models'' or ``about collaboration'' in the abstract; they are about domain-grounded deployment strategies under resource constraints.
